\documentclass[11pt]{article}
\usepackage[T1]{fontenc}
\usepackage{amsmath,amssymb,amsthm}
\usepackage{booktabs}
\usepackage[margin=1in]{geometry}
\usepackage{url}
\usepackage[hidelinks]{hyperref}

\newtheorem{theorem}{Theorem}
\newtheorem{lemma}[theorem]{Lemma}
\newtheorem{proposition}[theorem]{Proposition}
\theoremstyle{definition}
\newtheorem{definition}[theorem]{Definition}

\newcommand{\BRQ}{\mathrm{BR}_{\operatorname{int}(Q)}}
\newcommand{\supp}{\operatorname{supp}}

\title{Subtraction-Free Bounded Ratios in One Sign Chamber\\
at Degree Four in Three Variables}
\author{Xiaojian Zeng\\
PBC School of Finance, Tsinghua University, 100190\\
\texttt{xiaojianzeng@foxmail.com}}
\date{}
\hypersetup{
  pdftitle={Subtraction-Free Bounded Ratios in One Sign Chamber at Degree Four in Three Variables},
  pdfauthor={Xiaojian Zeng},
  pdfsubject={An exact Hilbert-basis proof for one closed sign chamber},
  pdfkeywords={Lorentzian polynomials, bounded ratios, subtraction-free, Hilbert basis, Normaliz}
}

\begin{document}
\maketitle

\begin{abstract}
For a product of four nonnegative ternary linear forms, Bathija, Rohatgi,
and Soskin conjectured that every integral point of the product
bounded-ratio cone gives a subtraction-free ratio.  We prove the conjectured
conclusion on one full-dimensional closed coordinate-sign chamber at
$(n,k)=(4,3)$.  In the coordinate order
\[
004,013,022,031,040,103,112,121,130,202,211,220,301,310,400,
\]
the chamber has weak sign pattern \(+++++--+++++-++\).  An exact
saturated-lattice computation gives a Hilbert basis of 134 elements, of
which 86 are not extreme-ray generators.  Exact polynomial-support
enumeration verifies the required containment for every Hilbert-basis
element, and sign compatibility propagates it to every lattice point in the
closed chamber.  The assertion for the other sign chambers, and hence the
original universal conjecture, remains unresolved.
\end{abstract}

\section{Introduction}

Lorentzian polynomials provide a common framework for log-concavity
phenomena in combinatorics and geometry \cite{BrandenHuh2020}.  Bathija,
Rohatgi, and Soskin introduced the corresponding bounded-ratio problem for
products of linear forms and formulated a coefficientwise
subtraction-free conjecture for all integral bounded ratios
\cite{BathijaRohatgiSoskin2026}.  Their Proposition~7.4 shows that at
$(n,k)=(4,3)$ the product bounded-ratio cone has 18 triangular extreme
rays.  Real conical generation by these rays does not, however, imply
generation of every lattice point by the primitive rays.

The degree-two predecessor is the rank-two Lorentzian-matrix conjecture of
Huang, Huh, Soskin, and Wang \cite[Conjecture~5.5]{HuangHuhSoskinWang2025}.
The higher-degree formulation uses the ordinary coefficients of the product
of linear forms and allows a constant depending on the exponent vector
\cite[Definitions~5.4 and~9.1]{BathijaRohatgiSoskin2026}.  We retain those
conventions throughout.

Our contribution is deliberately local: we prove the conjectured conclusion
for every integral point in one complete closed sign chamber.  The chamber
is full-dimensional in the 12-dimensional balanced space, and its Hilbert
basis has 86 non-ray elements.  Thus the result goes beyond checking the 18
source rays, but it makes no assertion about points with an incompatible
coordinate-sign pattern.

\section{Coefficients, ratios, and the chamber}

Let
\[
 H=H(4,3)=\{\alpha\in\mathbb Z_{\geq0}^{3}:|\alpha|=4\}
\]
and let $A=(a_{ri})$ be a formal $4\times3$ matrix.  Write
\begin{equation}\label{eq:product}
 Q_A(x_1,x_2,x_3)=\prod_{r=1}^{4}
 (a_{r1}x_1+a_{r2}x_2+a_{r3}x_3)
 =\sum_{\alpha\in H}P_\alpha(A)x^\alpha .
\end{equation}
These are ordinary coefficients:
\begin{equation}\label{eq:assignment}
 P_\alpha(A)=
 \sum_{\phi:\{1,2,3,4\}\to\{1,2,3\},\;|\phi^{-1}(i)|=\alpha_i}
 \prod_{r=1}^{4}a_{r,\phi(r)}.
\end{equation}
In particular, every monomial appearing in $P_\alpha$ has positive integer
coefficient.

For an integral exponent vector $\gamma=(\gamma_\alpha)_{\alpha\in H}$,
put
\[
 N_\gamma=\prod_{\gamma_\alpha>0}P_\alpha^{\gamma_\alpha},
 \qquad
 D_\gamma=\prod_{\gamma_\alpha<0}P_\alpha^{-\gamma_\alpha}.
\]
The ratio $N_\gamma/D_\gamma$ is subtraction-free if some
$C=C(\gamma)>0$ makes $CD_\gamma-N_\gamma$ coefficientwise nonnegative in
the twelve variables $a_{ri}$.

We use the fixed order
\begin{equation}\label{eq:order}
 (004,013,022,031,040,103,112,121,130,202,211,220,301,310,400)
\end{equation}
on $H$.  Let $K=\BRQ(4,3)$ and let
\[
 L=\left\{\gamma\in\mathbb Z^{H}:
       \sum_{\alpha\in H}\gamma_\alpha\alpha=0\right\}
\]
be the balanced lattice.  Every point of $K$ is balanced
\cite{BathijaRohatgiSoskin2026}.  Define the affine semigroup
\begin{equation}\label{eq:semigroup}
 S=K\cap L\cap
 \left\{\begin{array}{ll}
 \gamma_{103},\gamma_{112},\gamma_{301}\leq0,\\
 \gamma_\alpha\geq0&\text{for all other }\alpha\in H.
 \end{array}\right.
\end{equation}
Thus $S$ includes every boundary face and the zero vector.

\begin{theorem}\label{thm:main}
For every $\gamma\in S$, there is a constant $C(\gamma)>0$ such that
$C(\gamma)D_\gamma-N_\gamma$ is coefficientwise nonnegative in the
twelve parameters $a_{ri}$.
\end{theorem}

The theorem is exactly the subtraction-free assertion on the closed sign
chamber \(+++++--+++++-++\) in the order \eqref{eq:order}.  It does not
resolve the assertion for all of $K\cap\mathbb Z^H$.

\section{Two finite reductions}

For a polynomial with nonnegative coefficients, let $\supp(F)$ denote its
set of exponent vectors.

\begin{lemma}[Support criterion]\label{lem:support}
For a fixed integral $\gamma$, a positive subtraction-free constant exists
if and only if
\[
 \supp(N_\gamma)\subseteq\supp(D_\gamma).
\]
\end{lemma}

\begin{proof}
All nonzero coefficients of $N_\gamma$ and $D_\gamma$ are positive
integers.  If the containment holds, one may take
\[
 C=\max_{u\in\supp(N_\gamma)}
 \frac{[a^u]N_\gamma}{[a^u]D_\gamma}.
\]
Conversely, a monomial in the numerator support but not in the denominator
has a negative coefficient in $CD_\gamma-N_\gamma$ for every $C$.
\end{proof}

\begin{lemma}[Sign-compatible multiplication]\label{lem:multiply}
If $\gamma$ and $\eta$ both satisfy the weak signs in
\eqref{eq:semigroup}, then
\[
 N_{\gamma+\eta}=N_\gamma N_\eta,
 \qquad D_{\gamma+\eta}=D_\gamma D_\eta.
\]
Consequently, subtraction-freeness for two such vectors implies
subtraction-freeness for their sum.
\end{lemma}

\begin{proof}
At every coordinate, $\gamma_\alpha$ and $\eta_\alpha$ have the same weak
sign, so no numerator--denominator cancellation occurs.  If constants
$C_\gamma$ and $C_\eta$ work, then
\begin{align*}
 C_\gamma C_\eta D_\gamma D_\eta-N_\gamma N_\eta
 ={}&C_\eta D_\eta(C_\gamma D_\gamma-N_\gamma)\\
   &+N_\gamma(C_\eta D_\eta-N_\eta),
\end{align*}
which is coefficientwise nonnegative.
\end{proof}

It therefore suffices to find the Hilbert basis of $S$ and apply
Lemma~\ref{lem:support} to its elements.  The lattice in that computation
is essential: the primitive real extreme rays alone need not generate
$S$ as a monoid.

\section{Exact Hilbert-basis and support certificate}

For $\beta\in H(2,3)$ and $\{i,j,k\}=\{1,2,3\}$, let
$\rho_{\beta;i\mid jk}$ have positive coordinates at
$\beta+2e_i$ and $\beta+e_j+e_k$, negative coordinates at
$\beta+e_i+e_j$ and $\beta+e_i+e_k$, and zero coordinates elsewhere.
Proposition~7.4 of \cite{BathijaRohatgiSoskin2026} gives
\begin{equation}\label{eq:rays}
 K=\operatorname{cone}\{\rho_{\beta;i\mid jk}:
       \beta\in H(2,3),\ i\in\{1,2,3\}\},
\end{equation}
with 18 distinct extreme rays.

Let $B$ be the $3\times15$ matrix whose columns, in the order
\eqref{eq:order}, are the elements of $H$, and let $R$ be the
$15\times18$ matrix whose columns are the rays in \eqref{eq:rays}.  Exact
integer linear algebra gives
\[
 \operatorname{rank}B=3,\qquad \operatorname{rank}R=12,
 \qquad BR=0,
\]
and the twelve nonzero Smith invariants of $R$ are all $1$.  Hence the
integer span of the 18 rays is saturated in $\mathbb Z^{15}$ and, because
its rational span is $\ker B$, it equals $L=\ker(B)\cap\mathbb Z^{15}$.
This identifies the lattice used by the Hilbert-basis calculation with the
full balanced lattice rather than a finite-index sublattice.

In rank-12 lattice coordinates, the cone $K$ has 18 reconstructed support
halfspaces.  Intersecting them with the 15 weak coordinate-sign halfspaces
from \eqref{eq:semigroup} gives a pointed 12-dimensional cone.  Exact
double-description cuts give 48 primitive extreme rays.  The Hilbert basis
was then computed with Normaliz \cite{BrunsIchim2010,Normaliz} twice: once
from those 48 rays and once from all 33 defining halfspaces.  The two
calculations agree on the 15 irredundant support hyperplanes, the 48 extreme
rays, and a 134-element Hilbert basis.  Among the basis elements, 48 lie on
extreme rays and 86 do not.  If
$d_+(\gamma)=\sum_{\gamma_\alpha>0}\gamma_\alpha$, their distribution is
\begin{center}
\small
\begin{tabular}{c@{\quad}rrrrrrrrrr}
\toprule
$d_+$&2&3&4&5&6&7&8&9&10&12\\
\midrule
number&9&13&31&25&26&11&11&3&3&2\\
\bottomrule
\end{tabular}
\end{center}

Each returned vector was reconstructed in the original 15 coordinates and
checked against the three balance equations, the saturated lattice basis,
all 33 defining inequalities, and the sign conditions.  The polynomial
test begins directly from the assignment formula \eqref{eq:assignment}.
The support of a product is formed by exact Minkowski addition of finite
integer supports.

For efficiency, the verifier records in each row only the exponents of
$a_{r1}$ and $a_{r2}$.  This eight-coordinate compression loses no
information.  Indeed, balance implies
\(
 \sum_\alpha\gamma_\alpha=0
\)
because every $\alpha\in H$ has coordinate sum four.  Thus $N_\gamma$ and
$D_\gamma$ contain the same number $d_+(\gamma)$ of coefficient-polynomial
factors, counted with multiplicity.  In either product the total degree in
row $r$ is $d_+(\gamma)$, so the omitted exponent of $a_{r3}$ is forced by
the other two.

The exact support calculation gives
\begin{equation}\label{eq:all-supports}
 \supp(N_h)\subseteq\supp(D_h)
 \qquad\text{for every one of the 134 Hilbert-basis elements }h.
\end{equation}
There were no missing-support cases.  The largest numerator and denominator
supports encountered had 889 and 157,521 elements, respectively.  These
figures are bookkeeping checks; the proof-relevant output is the exact
containment \eqref{eq:all-supports} for every certified generator.

\begin{proof}[Proof of Theorem~\ref{thm:main}]
By the Hilbert-basis certificate, every nonzero $\gamma\in S$ is a
nonnegative integral sum of the 134 recorded elements.  By
\eqref{eq:all-supports} and Lemma~\ref{lem:support}, each summand is
subtraction-free.  All summands share the weak sign pattern, so repeated
application of Lemma~\ref{lem:multiply} proves subtraction-freeness of
$\gamma$.  For $\gamma=0$, one has $N_0=D_0=1$, and $C=1$ works.
\end{proof}

\section{Scope and reproducibility}

The strict coordinate hyperplanes cut the relative interior of $K$ into
4,074 realizable sign patterns.  The theorem treats the closure of one of
them.  The other 4,073 strict patterns have not been covered by this
Hilbert-basis argument, so Conjecture~9.2 of
\cite{BathijaRohatgiSoskin2026}, even at $(4,3)$, remains unresolved.  Nor
does this result assert that $C=1$ works for the 86 non-ray generators.

The ancillary program
\path{evidence/verify_chamber_hilbert.py} and its serialized certificate
\path{evidence/chamber-hilbert-basis.json} reproduce the integer-lattice,
cone, and support calculations.  The recorded run used Python~3, SymPy for
exact integer and rational linear algebra \cite{MeurerEtAl2017}, and
PyNormaliz~2.24 as the interface to Normaliz.  No floating-point decision is
used.  A separate audit reconstructed all monomial supports in the full
twelve-variable coordinates and recovered the same 134 exponent vectors
from an ambient-coordinate cone calculation.

The use of replayable exact verifiers follows the same methodological
principle as other computer-assisted inequality work in which every
machine-assisted assertion is paired with an exact rational verifier
\cite[Section~14]{Zeng2026Lens}.  That paper supplies methodological context
only; no result from it is used here.

Generative-AI systems assisted with literature screening, program
development and checking, and preparation of this manuscript.  The
mathematical evidence for Theorem~\ref{thm:main} is the exact finite
certificate and the proof reductions given above; generated text or model
judgments are not used as evidence.

\bibliographystyle{alpha}
\bibliography{references}

\end{document}
